% ============ 第2章 §2.3 Lebesgue 空间 (原书 pp.47-51) ============

\section{Lebesgue 空间}\label{sec:2.3}

不仅限于通常的绝对可积, 考虑种种可积性是有意义的. 例如, $L^{2}$ 空间具有内积这一运算, 由此得到在 $L^{1}$ 空间中不成立的反射性等. 这种性质得以成立的背景, 是由范数得到的开球的几何性质. 本书除 $L^{2}$ 之外不深入讨论这些. 又, 从调和分析学的立场, 仅考虑 $L^{1}$ 空间不够充分的迫切理由, 是 $L^{1}$ 中许多算子不连续, 而在 $L^{p}$ 空间 ($1 < p < \infty$) 中连续这一事实. 这里, 先列举考察 $L^{p}$ 空间所必需的积分不等式, 再定义 $L^{p}$ 空间并复习其性质.

\subsection*{积分不等式}

已述在调和分析学中, 一边大量使用积分一边进行种种分析, 但能精确求出积分值的情形几乎没有. 多数情形要看该积分值小于哪个数. 虽然无法求出值, 却要评价该积分的大小, 如本书所见, 这是项艰巨的作业. 作为基础的下述不等式是自然的, 但精确使用由它派生出的不等式是关键.

\begin{theorem}\label{thm:2.15}
设 $(X, \mathcal{B}, \mu)$ 为测度空间. 此时, 对可测函数 $f, g: X \to [0, \infty]$, 或可积函数 $f, g: X \to \mathbb{R}$, 若 $f(x) \leq g(x)$, 则积分的单调性
\[
\int_{X}f(x)\dd\mu(x) \leq \int_{X}g(x)\dd\mu(x)
\]
成立.
\end{theorem}

\begin{proof}
这里不作证明. 请参考适当的文献.
\end{proof}

由这个不等式得到的一连串基本不等式列举如下. 它们各自有名字, 本书中经常使用.

\begin{theorem}\label{thm:2.16}
设 $(X, \mathcal{B}, \mu)$ 为测度空间. 此时, 对可积函数 $f: X \to \mathbb{C}$, 积分的三角不等式
\[
\left|\int_{X}f(x)\dd\mu(x)\right| \leq \|f\|_{L^{1}(\mu)} = \int_{X}|f(x)|\dd\mu(x)
\]
成立.
\end{theorem}

\begin{proof}
这里不作证明. 请参考适当的文献.
\end{proof}

\begin{theorem}\label{thm:2.17}
设 $(X, \mathcal{B}, \mu)$ 为测度空间. 此时, 设 $f$ 为 $\mu$-可测函数且 $\lambda > 0$. 此时, Chebyshev 不等式
\begin{equation}\tag{2.6}
\mu\{|f| > \lambda\} \leq \frac{1}{\lambda}\int_{X}|f(x)|\dd\mu(x)
\label{eq:2.6}
\end{equation}
成立.
\end{theorem}

\begin{proof}
这里不作证明. 请参考适当的文献.
\end{proof}

\subsection*{$L^{p}$-范数}

来定义 $L^{p}$-范数. 设 $1 \leq p < \infty$. 对可测函数 $f: X \to \mathbb{C}$ 或 $\overline{\mathbb{R}}$,
\[
\|f\|_{L^{p}} = \|f\|_{L^{p}(\mu)} = \left(\int_{X}|f(x)|^{p}\dd\mu(x)\right)^{\frac{1}{p}}
\]
成立. $\|f\|_{L^{p}}$, $\|f\|_{L^{p}(\mu)}$ 称为 $f$ 的 $L^{p}$-范数. 用 $L^{p}(\mu)$ 表示这个量为有限的 Lebesgue 可测函数全体.

对 Lebesgue 可测函数 $f: X \to \mathbb{C}$ 或 $\overline{\mathbb{R}}$,
\[
\|f\|_{L^{\infty}(\mu)} = \inf\{\lambda > 0 : |f(x)| \leq \lambda \text{ a.e.}\}
\]
成立. $\|f\|_{L^{\infty}}$, $\|f\|_{L^{\infty}(\mu)}$ 称为 $f$ 的 $L^{\infty}(\mu)$-范数, 或略称为 $L^{\infty}$-范数. 用 $L^{\infty}(\mu)$ 表示这个量为有限的 Lebesgue 可测函数全体.

在 $L^{p}(\mu)$ 的定义中, 当 $(X, \mathcal{M}, \mu) = (\mathbb{R}^{n}, \mathcal{L}, \dd x)$ 时, 把 $L^{p}(\mu)$ 写作 $L^{p}(\mathbb{R}^{n})$, 用 $\|f\|_{L^{p}}$ 表示其范数.

\subsection*{Minkowski 不等式}

下面复习本书中经常使用的 Minkowski 不等式.

\begin{theorem}\label{thm:2.18}
设 $1 \leq p < \infty$. 设 $f$, $g$ 为可测集合 $X$ 上的可测函数. 此时, Minkowski 不等式 $\|f + g\|_{L^{p}(\mu)} \leq \|f\|_{L^{p}(\mu)} + \|g\|_{L^{p}(\mu)}$ 成立.
\end{theorem}

\begin{proof}
这里不作证明. 请参考适当的文献. 不过, 指出定理~\ref{thm:1.1} 的证明被变换到积分的文脉, 成为定理~\ref{thm:2.18} 的证明.
\end{proof}

\subsection*{Hölder 不等式}

下面复习本书中经常使用的 Hölder 不等式.

\begin{theorem}\label{thm:2.19}
设 $1 \leq p < \infty$, $q$ 为 $p$ 的调和共轭. 设 $f$, $g$ 为测度空间 $(X, \mathcal{B}, \mu)$ 上的可测函数.

(1) Hölder 不等式 $\|fg\|_{L^{1}(\mu)} \leq \|f\|_{L^{p}(\mu)}\|g\|_{L^{q}(\mu)}$ 成立. 不过, 在右边需要时, 请按 $0 \times \infty = \infty \times 0 = 0$ 解释.

(2) 【$L^{p}(\mu)$ 空间的对偶性】 $\|f\|_{L^{p}(\mu)} = \max\{\|f \cdot g\|_{L^{1}(\mu)} : \|g\|_{L^{q}(\mu)} \leq 1\}$ 成立.

(3) 【$L^{\infty}(\mu)$ 空间的对偶性】 $\|f\|_{L^{\infty}(\mu)} = \sup\{\|f \cdot g\|_{L^{1}(\mu)} : \|g\|_{L^{1}(\mu)} \leq 1\}$ 成立.
\end{theorem}

\begin{proof}
这里不作证明. 请参考适当的文献. 不过, 指出定理~\ref{thm:1.2} 及其直后相关命题的证明被变换到积分的文脉, 成为定理~\ref{thm:2.19} 的证明.
\end{proof}

\subsection*{卷积}

卷积在积分算子中也属于基本的, 出现在许多数学问题中. 偏微分方程的解算子多数由卷积得到, 这是显示卷积重要性的典型例子. 又, 分析数据时, 有通过取前后 $7$ 个数据的平均来去除数据间噪声的方法, 可以说这是在 $\mathbb{Z}$ 等离散空间上的卷积.

数字虽然没有精确的可信度, 数字也随时在变动, 但想说调和分析的 $45\%$ 左右的算子与卷积相关联的心情是有的.

设 $f$, $g$ 为 $\mathbb{R}^{n}$ 上的适当可测函数, $f$, $g$ 的卷积 $f*g$ 由
\[
f * g(x) = \int_{\mathbb{R}^{n}}f(x - y)g(y)\dd y
\]
定义. 一般, 当
\[
\int_{\mathbb{R}^{n}}|f(x - y)g(y)|\dd y = \infty
\]
时 $f * g(x)$ 没有意义, 但暂且只要有意义就考虑 $f * g(x)$. 本书特别是
\[
\int_{\mathbb{R}^{n}}|f(x - y)g(y)|\dd y = \infty
\]
的情形也多有考虑. 即使在这种情形, 在 $\pm$ 相互抵消的状况下也能考虑卷积, 许多情况下这极为有用, 在数学的其他领域带来重要结果. 设 $n \geq 3$, 偏微分方程之一的 Laplace 方程 $-\Delta u = f$ 的解具有
\[
u(x) = c_{n}\int_{\mathbb{R}^{n}}\frac{1}{|x - y|^{n-2}}f(y)\dd y \quad (x \in \mathbb{R}^{n})
\]
的形式, 所以可以看到卷积的形式. Schr\"odinger 方程 $\partial_{t}u - i\Delta u = 0$ 的解在初始值 $u_{0}$ 之下为
\[
u(x, t) = \frac{1}{\sqrt{(4\pi i t)^{n}}}\int_{\mathbb{R}^{n}}u_{0}(y)\exp\left(-\frac{|x - y|^{2}}{4it}\right)\dd y \quad (x \in \mathbb{R}^{n}, t > 0).
\]

卷积也出现在数学的其他领域. 设 $X$, $Y$ 为具有概率密度 $f$, $g$ 的独立随机变量. 因为 $X$ 具有概率密度 $f$,
\[
P(X > \lambda) = X \text{ 取大于 } \lambda \text{ 的值的概率} = \int_{\lambda}^{\infty}f(t)\dd t
\]
成立. 对 $Y$ 与 $g$ 也有对应的关系式成立. 此时, $X + Y$ 的概率密度由 $f * g$ 给出.

在上述情形中 $f * g(x)$ 没有意义时, 形式上令 $f * g(x) = 0$ 等, $f * g$ 暂且视为处处有定义.

\begin{theorem}[Young 不等式]\label{thm:2.20}
当 $1 \leq p, q, r \leq \infty$ 满足 $\frac{1}{r} = \frac{1}{p} + \frac{1}{q} - 1$ 时, 对 $f \in L^{p}(\mathbb{R}^{n})$, $g \in L^{q}(\mathbb{R}^{n})$, $f * g$ 有意义, 且 $\|f * g\|_{L^{r}} \leq \|f\|_{L^{p}}\|g\|_{L^{q}}$ 成立.
\end{theorem}

\begin{proof}
这里不作证明. 请参考适当的文献.
\end{proof}

\begin{exercise}\label{exer:2.1}
对 $A, B \geq 0$ 与 $0 < p < \infty$, 通过考虑 $A \leq B$ 与 $B \leq A$ 两种情形, 证明 $2^{-1}(A^{p} + B^{p}) \leq (A + B)^{p} \leq 2^{p}(A^{p} + B^{p})$. 【注意】由于具有对称性, 可假定前者来考虑.
\end{exercise}

\begin{exercise}\label{exer:2.2}
连续函数 $f: (0, \infty) \to \mathbb{R}$ 的 $L^{p}$-范数由
\[
\|f\|_{L^{p}(0, \infty)} = \left(\int_{0}^{\infty}|f(t)|^{p}\dd t\right)^{\frac{1}{p}}
\]
给出. 设 $a$ 为实数, 回答以下问题.

(1) 设 $f(x) = (1 + x)^{-1}$. 计算 $f$ 的 $L^{1}(0, \infty)$-范数.

(2) 设 $-\infty < a < 1$. 设 $f(x) = (1 + x)^{-a}$. 计算 $f$ 的 $L^{1}(0, \infty)$-范数.

(3) 设 $a > 1$. 设 $f(x) = (1 + x)^{-a}$. 计算 $f$ 的 $L^{1}(0, \infty)$-范数.

(4) 设 $f(x) = (1 + x)^{-a}$. $f$ 属于 $L^{1}(0, \infty)$ 的必要充分条件是什么?

(5) 设 $f(x) = (1 + x)^{-a}$. $f$ 属于 $L^{2}(0, \infty)$ 的必要充分条件是什么?
\end{exercise}

\begin{exercise}\label{exer:2.3}
令 $B(1) = \{|x| < 1\}$. 写出 $L^{p}(\mathbb{R}^{n})$ 的定义, 并证明当 $a < -\frac{n}{p}$ 时 $f(x) = \chi_{\mathbb{R}^{n}\setminus B(1)}(x)|x|^{a}$ 属于 $L^{p}(\mathbb{R}^{n})$.
\end{exercise}
